\subsection{Quantitative Transfer Evaluation of the Framework}
\label{sec:framework_transfer}

We supplement the scenario comparisons with a quantitative test of whether
an equal-RTT calibration transfers to different capacities and delays.
This analysis was designed during revision after the archived results were
available. Calibration and evaluation runs are disjoint, but the test is
retrospective; its design and data split are documented in the supplement.
The analysis script retains all inputs and reports the constant-baseline
comparison, including the contribution of queue conditioning.

Calibration uses only the 15 BBRv3 E1 runs with equal 20-ms RTTs at
100 Mbps, five at each queue target $q\in\{0.5,2,10\}$. At each target,
we estimate mean $u$ and mean $J$. For a target queue ratio between knots,
we linearly interpolate the two means in $\log q$; outside the range,
we extend the adjacent segment and restrict predictions to the metric's
feasible interval. The comparator predicts the overall calibration mean,
independent of queue size. Neither model fits an E3 outcome. Capacity and
RTT effects cannot be estimated from this single-capacity, equal-RTT
calibration; transferring normalized rates therefore tests an explicit
assumption that the evaluation tests.

The primary transfer comparison uses the four 200-Mbps E3 cases, where
physical queue ratios agree with the frozen input labels: RTTs of 10 or
60 ms and $q\in\{1,4\}$. The four 50-Mbps cases are a separate comparison
using their corrected physical ratios of 4 and 16. Errors are computed
against each condition's five-run mean, for 20 BBRv3 runs in each group.
The 40 CUBIC runs remain a performance comparator and are not calibration
data for the BBRv3 predictor.

\begin{table}[htbp]
\centering
\small
\caption{Retrospective transfer errors for BBRv3 condition means. MAE is the
mean absolute error; normalized-goodput errors are fractions of capacity.
Queue interpolation uses E1 only. The constant predictor uses the same
calibration data. Each row evaluates four E3 conditions.}
\label{tab:framework_transfer}
\begin{tabular}{llrrr}
\toprule
Evaluation group & Metric & Queue MAE & Constant MAE & Queue max. error \\
\midrule
200 Mbps, actual $q=1,4$ & $u$ & 0.00264 & 0.00264 & 0.00400 \\
 & $J$ & 0.01010 & 0.01134 & 0.01464 \\
50 Mbps, actual $q=4,16$ & $u$ & 0.00309 & 0.00309 & 0.00366 \\
 & $J$ & 0.02027 & 0.01444 & 0.02499 \\
\bottomrule
\end{tabular}
\end{table}

Table~\ref{tab:framework_transfer} separates modest transfer error from
predictive added value. In the four input-matched cases, the queue-conditioned
Jain estimate has slightly smaller MAE than the constant predictor, but
normalized-goodput MAE is the same to the displayed precision. In the corrected 50-Mbps cases, Jain MAE is 0.02027 for queue
interpolation and 0.01444 for the constant estimate. This comparison
identifies the dependence of transfer estimates on their conditioning
variables. The largest Jain error is about
0.025; reporting only high aggregate utilization would conceal this
assessment limitation. These results quantify the framework's transfer
boundary without claiming that E3 validates the original queue thresholds.

We also instantiate the analytical assessment in
Section~\ref{sec:analytical_assessment} with illustrative requirements
$u_0=0.90$ and $j_0=0.95$. These values were selected for retrospective
illustration, not frozen as experimental acceptance thresholds. All 40
BBRv3 and all 40 CUBIC E3 runs meet the utilization requirement, while
40 BBRv3 and 27 CUBIC runs meet both requirements. The result is restricted
to equal-RTT same-algorithm competition; E2 provides the separate coexistence
assessment. The supplement reports
run-level decisions and a threshold sensitivity grid. Receiver-goodput and early-ping windows are reported separately. Sender
retransmission totals and queue snapshots provide complementary indicators
for the documented operating conditions.

\begin{table}[htbp]
\centering
\small
\caption{Sensitivity of the retrospective E3 rate/fairness assessment to
illustrative service requirements. Cells count runs meeting both limits,
from 40 runs per algorithm. They are not independent condition counts.}
\label{tab:assessment_requirements}
\begin{tabular}{rrrr}
\toprule
Minimum $u$ & Minimum $J$ & BBRv3 & CUBIC \\
\midrule
0.90 & 0.95 & 40/40 & 27/40 \\
0.95 & 0.95 & 0/40 & 27/40 \\
0.90 & 0.99 & 38/40 & 19/40 \\
\bottomrule
\end{tabular}
\end{table}

The higher utilization requirement in Table~\ref{tab:assessment_requirements}
reverses which algorithm meets the joint requirement more often: no
BBRv3 run reaches $u=0.95$, whereas all CUBIC runs do. Tightening fairness
instead reduces both algorithms' pass counts. This measured throughput--allocation tradeoff makes the service
preference explicit under equal RTTs. The full
threshold grid and per-condition variability are provided in the supplement.

The frozen F07 predictions are retained separately from this new analysis.
They specify qualitative ``near-capacity'' or ``near-one'' outcomes without
numerical acceptance margins, and request medians rather than the reported
full-run means. Four cases also have inconsistent physical queue inputs.
The new transfer test reports numerical error directly; the historical
prediction sheet remains in the supplement as a separate record.
